\documentclass[11pt]{article}
\usepackage[margin=1in]{geometry}
\usepackage{amsmath,amssymb,amsthm}
\usepackage{graphicx}
\usepackage[hidelinks]{hyperref}

\newtheorem{theorem}{Theorem}
\newtheorem{corollary}{Corollary}
\newtheorem{remark}{Remark}

\newcommand{\C}{\mathbb C}
\newcommand{\eps}{\varepsilon}
\newcommand{\ran}{\operatorname{ran}}
\newcommand{\spanop}{\operatorname{span}}

\title{On the Operator Alpha-Ratio and Alpha-Reiter Problems}
\author{Run \texttt{fa\_banach\_001}, lane 7}
\date{June 20, 2026}

\begin{document}
\maketitle

\section*{Source Problems}

The source paper is N. Monod, \emph{Fixed points in convex cones},
arXiv:1701.05537.  Near the end of the paper, after discussing
$\ell^p$-variants of the $f$-ratio and $f$-Reiter properties, the author
introduces the following Hilbert-space analogues.  A unitary representation
$\pi$ of a group $G$ on a Hilbert space $V$ has the operator $\alpha$-ratio
property, for a continuous operator $\alpha\neq 0$, if for every finite
$S\subseteq G$ and every $\eps>0$ there is $v\in V$ such that
\[
  \|\alpha\pi(s)v\| < (1+\eps)\|\alpha v\|\qquad (s\in S).
\]
It has the stronger operator $\alpha$-Reiter property if one can choose $v$
such that
\[
  \|\alpha(\pi(s)v-v)\| < \eps\|\alpha v\|\qquad (s\in S).
\]
Problems 47 and 48 ask to clarify these properties, especially for projections
$\alpha$.

\begin{center}
\includegraphics[width=\textwidth]{figures/open_problem_crop.png}
\end{center}

The following theorem gives a direct clarification.  The ratio property is
automatic as stated; the Reiter property is exactly the absence of a relative
spectral gap for the seminorm induced by $\alpha$.

\section*{Idea of the Proof}

The ratio condition only asks that the $\alpha$-size of finitely many unitary
translates not exceed the original $\alpha$-size by much.  A vector which
almost attains the operator norm of $\alpha$ already has this property, because
unitarity gives $\|\pi(s)v\|=\|v\|$ and hence
$\|\alpha\pi(s)v\|\leq \|\alpha\|\,\|v\|$.

For Reiter, there is no such norm-attainment trick.  The obstruction is exactly
a Poincare inequality measured by $\alpha$: a finite set $S$ prevents
$\alpha$-Reiter vectors precisely when the $\alpha$-movement under $S$ is
bounded below by the $\alpha$-size of the vector.

\begin{theorem}
\label{thm:main}
Let $\pi:G\to \mathcal U(V)$ be a unitary representation on a Hilbert space
and let $0\neq \alpha\in B(V)$.

\begin{enumerate}
\item The representation has the operator $\alpha$-ratio property.

\item The representation has the operator $\alpha$-Reiter property if and only
if, for every finite $S\subseteq G$,
\[
  \inf_{\alpha v\neq 0}
  \frac{\sum_{s\in S}\|\alpha(\pi(s)v-v)\|^2}{\|\alpha v\|^2}=0.
\]
Equivalently, for no finite $S\subseteq G$ does there exist a constant
$c>0$ such that
\[
  \sum_{s\in S}\|\alpha(\pi(s)v-v)\|^2
  \geq c\|\alpha v\|^2
  \qquad\text{for all }v\in V.
\]
\item If $\alpha$ is bounded below, then the operator $\alpha$-Reiter property
is equivalent to the ordinary existence of almost invariant vectors for
$\pi$.
\end{enumerate}
\end{theorem}

\begin{proof}
Fix a finite set $S\subseteq G$ and $\eps>0$.  Since $\alpha\neq 0$, choose a
unit vector $v$ with
\[
  \|\alpha v\|>\frac{\|\alpha\|}{1+\eps}.
\]
For every $s\in S$, unitarity gives
\[
  \|\alpha\pi(s)v\|\leq \|\alpha\| < (1+\eps)\|\alpha v\|.
\]
This proves the operator $\alpha$-ratio property.

For the Reiter assertion, the displayed infimum is zero exactly when, for
every $\delta>0$, there is $v$ with $\alpha v\neq 0$ and
\[
  \sum_{s\in S}\|\alpha(\pi(s)v-v)\|^2 < \delta^2\|\alpha v\|^2.
\]
This implies the defining Reiter inequalities for all $s\in S$ after choosing
$\delta=\eps$.  Conversely, the defining inequalities imply the displayed
sum estimate with $\delta^2=|S|\eps^2$.  Thus the infimum formulation and the
original definition are equivalent.  The second formulation is just the
negation of the existence of a positive lower bound for the same Rayleigh
quotient.

Finally suppose that $\alpha$ is bounded below: there are constants
$0<m\leq M<\infty$ such that
\[
  m\|w\|\leq \|\alpha w\|\leq M\|w\|\qquad(w\in V).
\]
If $\pi$ has almost invariant vectors, then for a given finite $S$ choose
$v\neq 0$ with
\[
  \|\pi(s)v-v\|<\eps m M^{-1}\|v\|\qquad(s\in S).
\]
It follows that
\[
  \|\alpha(\pi(s)v-v)\|
  \leq M\|\pi(s)v-v\|
  < \eps m\|v\|
  \leq \eps\|\alpha v\|.
\]
Thus $\alpha$-Reiter holds.  Conversely, if $\alpha$-Reiter holds, choose
$v$ with
\[
  \|\alpha(\pi(s)v-v)\|<\eps m M^{-1}\|\alpha v\|
  \qquad(s\in S).
\]
Then
\[
  m\|\pi(s)v-v\|\leq \|\alpha(\pi(s)v-v)\|
  < \eps m M^{-1}\|\alpha v\|
  \leq \eps m\|v\|,
\]
so $\|\pi(s)v-v\|<\eps\|v\|$ for all $s\in S$.  Hence $\pi$ has almost
invariant vectors.
\end{proof}

It is sometimes useful to write the criterion in operator form.  If
\[
  A_S=\sum_{s\in S}(\pi(s)-I)^*\alpha^*\alpha(\pi(s)-I),
  \qquad B=\alpha^*\alpha,
\]
then $\alpha$-Reiter fails exactly when $A_S\geq cB$ for some finite $S$ and
some $c>0$.

\begin{corollary}[Finite-rank projections]
\label{cor:finite-rank}
Let $P$ be an orthogonal projection with finite-dimensional range $W$.  Then
$\pi$ has the operator $P$-Reiter property if and only if, for every finite
$S\subseteq G$, there exists $v\in V$ such that
\[
  Pv\neq 0
  \quad\text{and}\quad
  P(\pi(s)v-v)=0\qquad(s\in S).
\]
\end{corollary}

\begin{proof}
For fixed $S$, consider the linear map
\[
  L:V\longrightarrow W\oplus W^S,\qquad
  L(v)=\big(Pv,\ (P(\pi(s)v-v))_{s\in S}\big).
\]
The range of $L$ is a subspace of the finite-dimensional space
$W\oplus W^S$, hence is closed.

If the exact condition holds, then the $P$-Reiter inequalities hold with
zero left hand side for this finite $S$.  Conversely, if $P$-Reiter holds,
choose $v_n$ with $\|Pv_n\|=1$ and
\[
  \sum_{s\in S}\|P(\pi(s)v_n-v_n)\|^2\longrightarrow 0.
\]
After passing to a subsequence, compactness of the unit sphere of $W$ gives
$Pv_n\to w$ with $\|w\|=1$.  Hence $(w,0)$ belongs to the closure of
$L(V)$, and therefore to $L(V)$ itself.  This gives a vector $v$ satisfying
the exact condition.
\end{proof}

\begin{corollary}[Rank-one projections]
Let $P_\xi$ be the orthogonal projection onto the line spanned by a non-zero
vector $\xi$.  Then $\pi$ has the operator $P_\xi$-Reiter property if and only
if, for every finite $S\subseteq G$,
\[
  \xi\notin
  \spanop\{\pi(s)^{-1}\xi-\xi:\ s\in S\}.
\]
\end{corollary}

\begin{proof}
By Corollary~\ref{cor:finite-rank}, for each finite $S$ we need a vector
$v$ such that $\langle v,\xi\rangle\neq 0$ and
\[
  \langle \pi(s)v-v,\xi\rangle=0\qquad(s\in S).
\]
The latter equations say exactly that $v$ annihilates the finite-dimensional
space
\[
  M_S=\spanop\{\pi(s)^{-1}\xi-\xi:\ s\in S\}.
\]
Such a $v$ with $\langle v,\xi\rangle\neq0$ exists precisely when
$\xi\notin M_S$.
\end{proof}

\section*{Consequences and Limitations}

The operator $\alpha$-ratio property, as formulated in the source problem, is
therefore vacuous for bounded non-zero operators: it holds for every unitary
representation of every group.

The operator $\alpha$-Reiter property is genuinely representation-theoretic.
For $\alpha=I$ it is exactly the usual almost-invariant-vector condition.  In
particular, for the left regular representation it is equivalent to
amenability of the group.  More generally, when $\alpha$ is bounded below,
Theorem~\ref{thm:main} shows that no new condition is obtained beyond ordinary
almost invariance.

For projections with finite-dimensional range, the approximate condition
collapses to the exact finite-dimensional condition in
Corollary~\ref{cor:finite-rank}.  This gives a concrete answer to the
finite-rank projection subcase of Problem 48.  The packet does not claim to
classify all infinite-rank projections, nor does it address the separate open
problems about products, supramenability, or finitely generated groups of
exponential growth with the fixed-point property for cones.

\section*{Novelty Check}

The run indexes were searched for arXiv:1701.05537 and for phrases including
\texttt{operator alpha-ratio}, \texttt{operator alpha-Reiter},
\texttt{fixed points in convex cones}, \texttt{conical fixed-point property},
and \texttt{supramenability}.  No existing packet for these endpoint operator
questions was found.  Web searches for exact phrases such as
\texttt{"operator alpha-ratio"}, \texttt{"operator alpha-Reiter"}, and
variants including \texttt{"Fixed points in convex cones"} did not return a
later explicit answer to Problems 47 or 48.  The result here is elementary,
but it appears to clarify a point left open in the source problem statement.

\begin{thebibliography}{9}

\bibitem{Monod}
N. Monod,
\emph{Fixed points in convex cones},
arXiv:1701.05537, 2017.

\end{thebibliography}

\end{document}
